\subsection{合卷：征服的所得与未结的账}

清初的所得，是把军队占领、皇帝驻跸、部院文书和部分地方合作接成一个可以持续运作的朝廷。北京提供了旧帝国的行政中心；撤藩削除了将领家族长期据有军政的危险；海峡、蒙古和西南的经营又扩大了清廷能够动员资源的范围。

这些收益带有不能消失的成本。降官、旗人、降军、民夫、移民、盟旗与海疆居民并不承担同一种征发，也不在同一时刻接受同一种秩序。中央以册封、驻防、海禁、赈济和设治解决的，常是下一轮牧地、饷源、迁徙或身份争议的开始。帝国的边界因而不是一条完成线，而是无数地方关系被反复接入的结果。

\subsection{制度得失}

\textbf{所得：} 摄政、部院接收与旗军调度使新王朝能够将一次军事突破转化为跨区域的行政架构；撤藩削弱了独立军政中心。

\textbf{代价：} 财政必须长期供养驻军、转运与身份体系；地方社会承担战后复垦、移民、海禁和征发的差异性风险。

\textbf{留下的压力：} 多语边地与海峡需要的不是单一命令，而是持续的中介、补给和谈判。这些压力将贯穿所谓“盛世”的扩张。
